\section{Design Overview}
\label{sec:overview}

\Ansor is an automated tensor program generation framework.
\autoref{fig:system-overview} shows the overall architecture of \Ansor. The input of \Ansor is a set of to be optimized DNNs.
Ansor uses the operator fusion algorithm from Relay\cite{roesch2019relay} to convert DNNs from popular model formats (\textit{e.g.}, ONNX~\cite{bai2019onnx}, TensorFlow PB) to partitioned small subgraphs.
\Ansor then generates tensor programs for these subgraphs.
\Ansor has three major components: (1) a program sampler that constructs a large search space and samples diverse programs from it; (2) a performance tuner that fine-tunes the performance of sampled programs; (3) a task scheduler that allocates time resources for optimizing multiple subgraphs in the DNNs.

\textbf{Program sampler.} One key challenge \Ansor has to address is generating a large search space for a given computational graph.
To cover diverse tensor programs with various high-level structures and low-level details, \Ansor utilizes a hierarchical representation of the search space with two levels: \textit{sketch} and \textit{annotation}   (\autoref{sec:program-sampling}). 
\Ansor defines the high-level structures of programs as sketches and leaves billions of low-level choices (\textit{e.g.}, tile size, parallel, unroll annotations) as annotations.
This representation allows \Ansor to enumerate high-level structures flexibly and sample low-level details efficiently.
\Ansor includes a program sampler that randomly samples programs from the space to provide comprehensive coverage of the search space. 

\textbf{Performance tuner.} The performance of randomly sampled programs is not necessarily good. The next challenge is to fine-tune them.
\Ansor employs evolutionary search and a learned cost model to perform fine-tuning iteratively (\autoref{sec:performance-fine-tuning}).
At each iteration, \Ansor uses re-sampled new programs as well as good programs from previous iterations as the initial population to start the evolutionary search.
Evolutionary search fine-tunes programs by mutation and crossover which perform out-of-order rewrite and address the limitation of sequential construction.
Querying the learned cost model is orders of magnitude faster than actual measurement, so we can evaluate thousands of programs in seconds.

\textbf{Task scheduler.} Using program sampling and performance fine-tuning allows \Ansor to find high-performance tensor programs for a computational graph.
Intuitively, treating a whole DNN as a single computational graph and generating a full tensor program for it could potentially achieve the optimal performance. This, however, is inefficient because it has to deal with the unnecessary exponential explosion of the search space.
Typically, the compiler partitions the large computational graph of a DNN into several small subgraphs~\cite{roesch2019relay, chen2018tvm}. This partition has a negligible effect on the performance thanks to the layer-by-layer construction nature of DNNs.
This brings the final challenge of \Ansor: how to allocate time resources when generating programs for multiple subgraphs.
The task scheduler  (\autoref{sec:task-scheduler}) in \Ansor uses a scheduling algorithm based on gradient descent to allocate resources to the subgraphs that are more likely to improve the end-to-end DNN performance.


\begin{figure}[t]
	\centering
	\includegraphics[width=0.9\columnwidth]{figures/overview.pdf}
	\vskip -0.8em
	\caption{System Overview.
		The gray arrows show the flow of extracting subgraphs from deep learning models and generating optimized programs for them.
		The green arrows mean the measurer returns profiling data to update the status of all components in the system.}
	\postcap 
	\label{fig:system-overview}
\end{figure}
